\chapter{Rachunek Kirby'ego: operacje na ramowanych splotach i grafach}
\label{ch:tg-kirby-full}

Rachunek Kirby'ego pozwala porównywać różne opisy
tej samej trójrozmaitości jako brzegu czterorozmaitości
uchwytowej. Ważne jest rozróżnienie dwóch pytań:
czy dwie operacje dają tę samą rozmaitość \emph{brzegową}
i czy zachowują całą \emph{czterorozmaitość}.
Przesunięcie uchwytu zachowuje obie;
rozdmuchanie i kontrakcja składnika o ramowaniu $\pm1$
mogą natomiast zmienić czterorozmaitość,
choć nie zmienia zorientowanej trójrozmaitości brzegowej.

\section{Ramowanie, równoległa kopia i przesunięcie uchwytu}
\label{sec:tg-kirby-slides}

\begin{definicja}[Przesunięcie uchwytu indeksu dwa]
Niech $L=K_1\sqcup\cdots\sqcup K_r\subset S^3$
będzie zorientowanym splotem o ustalonych całkowitych
ramowaniach. \emph{Przesunięciem uchwytu $K_i$ po $K_j$}
($i\ne j$) nazywamy zastąpienie $K_i$ sumą pasmową
z równoległą kopią $K_j'$ wyznaczoną ramowaniem $K_j$.
Pasmo wybieramy z wnętrzem rozłącznym od całego
splotu oraz od kopii równoległej. Końce pasma
przyłączamy do $K_i$ i $K_j'$. Nowa składowa
otrzymuje ramowanie indukowane przez sumę pasmową
ramowanych krzywych: również ich równoległe kopie
łączymy w kołnierzu pasma. Nie wybieramy nowego
ramowania niezależnie od tego ruchu.
Orientację kopii możemy wybrać zgodną lub przeciwną;
w obliczeniach algebraicznych odpowiada to znakom $+$ i $-$.
\end{definicja}

\begin{stwierdzenie}[Zmiana macierzy przecięcia]
\label{prop:tg-kirby-congruence}
Niech $W_L$ będzie czterorozmaitością złożoną z $B^4$
i $r$ uchwytów indeksu dwa, a $Q_L$ jej macierzą przecięcia.
Po przesunięciu uchwytu $i$ po uchwycie $j$ ze znakiem
$\varepsilon\in\{\pm1\}$ otrzymujemy macierz
$Q_{L'}=P^TQ_LP$, gdzie $P\in GL_r(\mathbb Z)$
zastępuje wektor bazowy $e_i$ przez
$e_i+\varepsilon e_j$ i pozostawia pozostałe bez zmian.
Ponadto $W_L$ i $W_{L'}$ są dyfeomorficzne
z zachowaniem orientacji, a indukowany dyfeomorfizm
ich brzegów również zachowuje orientację.
Nie wymaga się, aby był identycznością na brzegu.
\end{stwierdzenie}

\begin{proof}
Z geometrycznego twierdzenia o przesuwaniu uchwytów,
przyjętego w dowodzie stwierdzenia
\ref{st:kt-handle-slide}, otrzymujemy dyfeomorfizm
zachowujący orientację. Nie wyprowadzamy istnienia
tego dyfeomorfizmu z kongruencji macierzy.
W tym samym dowodzie wykazano, że względem
identyfikacji grup homologii nowa klasa rdzenia
jest równa $e_i+\varepsilon e_j$.

W konwencji kolumnowej $P=I+\varepsilon E_{ji}$,
gdzie $E_{ji}$ ma jedynkę w pozycji $(j,i)$
i zera w pozostałych pozycjach. Ponieważ $i\ne j$,
mamy $E_{ji}^2=0$, więc
$P^{-1}=I-\varepsilon E_{ji}$.
Ponadto $\det P=1$, gdyż dodanie wielokrotności
jednej kolumny macierzy jednostkowej do innej
nie zmienia jej wyznacznika. Biliniowość formy
przecięcia daje następujące wpisy:
\[
\begin{aligned}
q'_{ii}&=q_{ii}+2\varepsilon q_{ij}+q_{jj},\\
q'_{ik}&=q_{ik}+\varepsilon q_{jk}\quad(k\ne i),\\
q'_{kl}&=q_{kl}\quad(k,l\ne i).
\end{aligned}
\]
Są to dokładnie wpisy $P^TQ_LP$.
Zatem wyznacznik pozostaje niezmieniony,
a prawo bezwładności Sylvestera zapewnia
zachowanie sygnatury, również przy osobliwej
macierzy przecięcia.
\end{proof}

\begin{przyklad}[Dwie krzywe z liczbą splecenia $\ell$]
Niech $Q=\begin{psmallmatrix}a&\ell\\\ell&b\end{psmallmatrix}$.
Po dodaniu drugiego wektora do pierwszego otrzymujemy
\[
Q'=\begin{pmatrix}a+2\ell+b&\ell+b\\\ell+b&b\end{pmatrix}.
\]
Wpis diagonalny zwiększył się o $2\ell+b$,
a nowa liczba splecenia wynosi $\ell+b$.
Bezpośrednio
$\det Q'=(a+2\ell+b)b-(\ell+b)^2=ab-\ell^2$.
Rachunek ten jednocześnie kontroluje ramowanie,
liczby splecenia i niezmienniczość wyznacznika.
\end{przyklad}

\begin{lemat}[Reguła \emph{slam-dunk}]
\label{lem:tg-kirby-slam-dunk}
Niech $U\cup V\subset S^3$ będzie standardowym zorientowanym
splotem Hopfa, którego składowe mają liczbę splecenia $+1$.
Niech $r\in\mathbb Z$, $s\in\mathbb Q\setminus\{0\}$.
Wykonanie chirurgii o współczynniku całkowitym $r$ na $U$
i wymiernym $s$ na jego geometrycznym meridianie $V$
jest równoważne wykonaniu na pojedynczym węźle trywialnym
chirurgii o współczynniku $r-1/s$. Wszystkie
współczynniki odczytujemy względem preferowanych
długości i z ustaloną orientacją $S^3$.
Warunek $r\in\mathbb Z$ jest istotny:
formuła w tej postaci nie obowiązuje dla dwóch
dowolnych wymiernych ramowań.
\end{lemat}
\begin{proof}[Opis sklejenia torusów]
Zewnętrze splotu Hopfa jest dyfeomorficzne
z $T^2\times[0,1]$: w modelu
$S^3=\{(z,w)\in\mathbb C^2:|z|^2+|w|^2=1\}$
obie składowe są okręgami współrzędnościowymi.
Po usunięciu ich małych otoczeń współrzędna
$|z|$ przebiega domknięty przedział zawarty
w $(0,1)$, a argumenty $z,w$ są niezależnymi
współrzędnymi okręgowymi. Na jego dwóch torusach brzegowych
możemy wybrać meridiany $\mu_U,\mu_V$ i długości
$\lambda_U,\lambda_V$ tak, że
$[\lambda_U]=[\mu_V]$ i $[\lambda_V]=[\mu_U]$
po identyfikacji w przestrzeni produktu:
meridian jednej składowej biegnie w kierunku
argumentu zanikającej przy niej współrzędnej,
a preferowana długość w kierunku argumentu
niezanikającej współrzędnej. Znaki dobieramy
zgodnie z liczbą splecenia $+1$.
Napiszmy $s=c/d$ w postaci nieskracalnej,
gdzie $c\ne0$, $d>0$ oraz
$\gcd(c,d)=1$.
Najpierw wykonajmy wypełnienie torusa brzegowego $U$.
Ponieważ $r$ jest całkowite, meridian dołączanego
torusa pełnego ma nachylenie $r[\mu_U]+[\lambda_U]$.
Zewnętrze splotu jest produktem, więc po wypełnieniu
jednego jego końca otrzymujemy torus pełny $M$:
do dołączonego torusa pełnego dołożono jedynie kołnierz
jego brzegu. Na pozostałym torusie brzegowym $V$
meridian $M$ jest reprezentowany przez
$[\mu_V]+r[\lambda_V]$.

Utożsamiamy teraz $M$ z zewnętrzem węzła trywialnego.
W nowej bazie chirurgicznej na jego brzegu przyjmujemy
\[
 [\mu']=-[\lambda_V],\qquad
 [\lambda']=[\mu_V]+r[\lambda_V].
\]
Macierz tych dwóch wektorów w bazie
$([\mu_V],[\lambda_V])$ ma wyznacznik
$\det\begin{psmallmatrix}0&1\\-1&r\end{psmallmatrix}=1$.
Drugi wektor jest meridianem torusa pełnego $M$
(czyli preferowaną długością węzła trywialnego),
a pierwszy wraz z nim tworzy całkowitą bazę.
Identyfikacja z zewnętrzem węzła trywialnego
istnieje, ponieważ zachowujący orientację dyfeomorfizm
torusa, który przenosi meridian torusa pełnego na
meridian (z dokładnością do znaku), przedłuża się
do dyfeomorfizmu torusów pełnych. Istotnie w bazach meridian--długość macierz
zachowująca orientację i nachylenie meridianu
ma postać
$\begin{psmallmatrix}\delta&k\\0&\delta\end{psmallmatrix}$,
gdzie $\delta\in\{\pm1\}$ i $k\in\mathbb Z$.
Macierze z $\delta=1$ realizują dyfeomorfizmy
$(z,u)\mapsto(zu^k,u)$ torusa pełnego
$D^2\times S^1$; macierz $-I$ realizuje
$(z,u)\mapsto(\overline z,\overline u)$.
Oba typy odwzorowań zachowują orientację.
Każdą z zapisanych macierzy otrzymujemy przez
ich złożenie. Dla ogólnego dyfeomorfizmu brzegu
korzystamy ponadto z faktu, że dyfeomorfizm torusa
jest izotopijny do liniowego przedstawiciela
swojego działania na $H_1$; izotopię przedłużamy
w kołnierzu brzegu. Uzasadnia to niezależność
wypełnienia od wyboru drugiego wektora bazy
przy ustalonym nachyleniu meridianu.

Pozostałe wypełnienie, na $V$, narzuca krzywą
$c[\mu_V]+d[\lambda_V]$. Z powyższej zmiany bazy
wynika $[\lambda_V]=-[\mu']$ oraz
$[\mu_V]=[\lambda']+r[\mu']$, zatem
\[
 c[\mu_V]+d[\lambda_V]
 =(rc-d)[\mu']+c[\lambda'].
\]
Jest to nachylenie pierwotne, ponieważ
$\gcd(rc-d,c)=\gcd(d,c)=1$.
Otrzymujemy więc chirurgię na pojedynczym węźle
trywialnym o współczynniku $(rc-d)/c=r-1/s$,
także wtedy, gdy $rc-d=0$.
Wniosek dotyczy zorientowanej rozmaitości otrzymanej
przez sklejenie torusów, a nie tylko jej homologii.

Dla uzasadnienia konieczności warunku
całkowitości $r$ porównajmy dwa
rachunki pierwszej homologii. Przy
$r=1/2$ i $s=4$ dwie relacje
chirurgii na splocie Hopfa brzmią
$\mu_U+2\mu_V=0$ i
$\mu_U+4\mu_V=0$, więc
$H_1\cong\mathbb Z/2\mathbb Z$:
pierwsza relacja eliminuje $\mu_U=-2\mu_V$,
a druga staje się dokładnie $2\mu_V=0$.
Natomiast chirurgia na pojedynczym
węźle trywialnym ze współczynnikiem
$r-1/s=1/4$ ma
$H_1=0$. Rozmaitości te nie mogą
być dyfeomorficzne; niedozwolone jest
więc stosowanie opisanej eliminacji
przy dwóch dowolnych wymiernych
współczynnikach.
\end{proof}

Powyższy dowód dotyczy izolowanego standardowego
splotu Hopfa. Nie dowodzi sam przez się, że przy
usuwaniu składowej z większego splotu pozostałe
składowe oraz ich ramowania pozostają niezmienione.
W takim zastosowaniu trzeba kontrolować sklejenie
względem pozostałych części zewnętrza; przykład
takiej geometrycznej redukcji ramion podano
w dowodzie twierdzenia \ref{thm:tg-star-seifert}.

\begin{przyklad}[Początek redukcji łańcucha]
Dwa kolejne wierzchołki plumbingowe o wagach
$-5,-2$ odpowiadają splotowi Hopfa o tych samych
ramowaniach. Po usunięciu składowej o ramowaniu
$-2$ otrzymujemy współczynnik
$-5-1/(-2)=-9/2$. Jest to właśnie ułamek
$-[5,2]^-$, który pojawi się przy konfiguracji
$C_3$ w rozdziale o wymiernej kontrakcji.
\end{przyklad}

\section{Rozdmuchanie i kontrakcja składnika o ramowaniu \texorpdfstring{$\pm1$}{+/-1}}
\label{sec:tg-kirby-blowdown}

\begin{twierdzenie}[Usuwanie składowej o ramowaniu $\pm1$]
\label{thm:tg-kirby-blowdown}
Niech $L=U\sqcup K_1\sqcup\cdots\sqcup K_r$
będzie całkowicie ramowanym splotem w $S^3$,
przy czym $U$ jest węzłem trywialnym o ramowaniu
$\varepsilon\in\{+1,-1\}$.
Niech $\ell_i=\operatorname{lk}(U,K_i)$ i niech
$B$ będzie macierzą splecenia składowych $K_i$.
Wtedy macierz całego splotu ma postać
$Q=\begin{psmallmatrix}\varepsilon&\ell^T\\\ell&B\end{psmallmatrix}$.
Po kontrakcji $U$ i odpowiednim pełnym skręceniu
pozostałych pasm otrzymujemy splot $L'$ o macierzy
\[
B'=B-\varepsilon\,\ell\ell^T.
\]
Chirurgie na $L$ i $L'$ dają orientacyjnie dyfeomorficzne
trójrozmaitości brzegowe. Na poziomie form przecięcia
zachodzi całkowita kongruencja
$Q\cong[\varepsilon]\oplus B'$.
Odpowiadające czterorozmaitości złożone z kuli
i uchwytów indeksu dwa nie są dyfeomorficzne:
usunięcie składnika o ramowaniu $\pm1$
zmniejsza ich drugą liczbę Bettiego o jeden,
zgodnie z twierdzeniem
\ref{thm:kt-handle-intersection}.
\end{twierdzenie}

\begin{proof}
Dla generatorów homologii $u,e_1,\ldots,e_r$
formy przecięcia mamy $u^2=\varepsilon$,
$u\cdot e_i=\ell_i$, $e_i\cdot e_j=B_{ij}$.
Przyjmijmy
$f_i=e_i-\varepsilon\ell_i u$. Ponieważ
$\varepsilon^2=1$, otrzymujemy
$u\cdot f_i=\ell_i-\varepsilon\ell_i\varepsilon=0$.
Ponadto
\[
f_i\cdot f_j=B_{ij}
-\varepsilon\ell_i\ell_j
-\varepsilon\ell_i\ell_j
+\varepsilon^2\ell_i\ell_j\varepsilon
=B_{ij}-\varepsilon\ell_i\ell_j.
\]
Zmiana bazy jest całkowita i unimodularna,
a zatem $Q\cong[\varepsilon]\oplus B'$.

\emph{Realizacja geometryczna.}
Zewnętrze węzła trywialnego $U$ jest torusem
pełnym, którego meridianem jest preferowana
długość $\lambda_U$ węzła. Skręt wzdłuż dysku
meridianowego tego torusa realizuje na brzegu
przekształcenie
\[
\mu_U\longmapsto\mu_U-\varepsilon\lambda_U,
\qquad \lambda_U\longmapsto\lambda_U.
\]
Można to zobaczyć w produktowych współrzędnych
$D^2\times S^1$: odwzorowanie
$(z,u)\mapsto(zu^{-\varepsilon},u)$
zachowuje orientację, jest dyfeomorfizmem
i dodaje $-\varepsilon$ meridianów do długości
torusa pełnego, czyli do $\mu_U$.
Wysyła nachylenie $\varepsilon\mu_U+\lambda_U$
na $\varepsilon\mu_U$, ponieważ $\varepsilon^2=1$.
Przedłuża się więc na dołączany torus pełny,
jak w dowodzie lematu \ref{lem:tg-kirby-slam-dunk}.
Po wypełnieniu otrzymujemy $S^3$ z obrazami
pozostałych ramowanych składowych. Jest to pełny
skręt pasm wzdłuż dysku rozpiętego na $U$,
z wykładnikiem $-\varepsilon$ w opisanej konwencji.
Dalsze wypełnienia wykonujemy wzdłuż obrazów
pierwotnych nachyleń, więc wynikowa zorientowana
trójrozmaitość nie zmienia się.

\emph{Obliczenie rzeczywistych ramowań po skręcie.}
Niech $n_i=B_{ii}$ oraz niech $\mu_i,\lambda_i$
będą meridianem i preferowaną długością $K_i$.
W pierwszej homologii zewnętrza całego splotu
meridiany tworzą bazę, a długości spełniają
(por. sekcja \ref{sec:kt-chirurgia-link})
\[
[\lambda_U]=\sum_j\ell_j[\mu_j],
\quad
[\lambda_i]=\ell_i[\mu_U]+
\sum_{j\ne i}B_{ij}[\mu_j].
\]
Wypełnienie składowej $U$ narzuca relację
$\varepsilon[\mu_U]+[\lambda_U]=0$.
Ponieważ $\varepsilon=\pm1$, eliminujemy
$[\mu_U]=-\varepsilon\sum_j\ell_j[\mu_j]$.
Pozostałe meridiany są nadal całkowitą bazą
homologii zewnętrza nowego splotu. Podstawienie daje
\[
[\lambda_i]=-\varepsilon\ell_i^2[\mu_i]
+\sum_{j\ne i}(B_{ij}-\varepsilon\ell_i\ell_j)[\mu_j].
\]
Przez $\lambda_i$ oznaczamy tu obraz starej
długości. Nowa preferowana długość $\lambda'_i$
musi mieć zerowy współczynnik przy własnym
meridianie; wobec tego jest określona na torusie
brzegowym przez
$\lambda'_i=\lambda_i+\varepsilon\ell_i^2\mu_i$.
Współczynniki przy innych meridianach dają
nowe liczby splecenia $B_{ij}-\varepsilon\ell_i\ell_j$.
Stare nachylenie chirurgiczne ma w nowej bazie postać
\[
n_i\mu_i+\lambda_i
=(n_i-\varepsilon\ell_i^2)\mu_i+\lambda'_i.
\]
Nowe ramowanie wynosi zatem
$n_i-\varepsilon\ell_i^2$.
Otrzymaliśmy wszystkie wpisy macierzy $B'$
bez wnioskowania o geometrii z samej kongruencji.
To kończy dowód.
\end{proof}

\begin{przyklad}[Liść o wadze $-1$ i centralny wierzchołek]
Załóżmy, że $U$ jest liściem grafu plumbingowego,
ma wagę $-1$ i łączy się tylko z wierzchołkiem
o wadze $b$, z liczbą splecenia $1$.
Wtedy $\varepsilon=-1$ i $\ell=(1,0,\ldots,0)$.
Z twierdzenia wynika $b'=b+1$, a pozostałe wagi
nie zmieniają się. Jest to klasyczna kontrakcja
liścia o wadze $-1$. Dla liścia z wagą $+1$
otrzymujemy odpowiednio $b'=b-1$.
\end{przyklad}

\begin{przyklad}[Kontrakcja węzła o dwóch sąsiadach]
Rozpatrzmy łańcuch z trzema wierzchołkami i wagami
$a,-1,b$. Macierz przecięcia wynosi
\[
Q=\begin{pmatrix}a&1&0\\1&-1&1\\0&1&b\end{pmatrix}.
\]
Usunięcie środkowego wierzchołka daje macierz
$B'=\begin{psmallmatrix}a+1&1\\1&b+1\end{psmallmatrix}$.
Nowe składowe są splecione raz, choć pierwotnie
końcowe wierzchołki nie były sąsiadami.
W szczególności nie wystarcza usunięcie wierzchołka
i zmiana samych wag: trzeba również dodać nową
krawędź grafu. Wynika to z pozadiagonalnego składnika
$-(-1)\ell_1\ell_2=1$.
\end{przyklad}

\begin{center}
\begin{tikzpicture}[scale=.92,>=Stealth]
 \node[tgBlueVertex] (a) at (-3,0) {$a$};
 \node[tgRedVertex] (x) at (-1.55,0) {$-1$};
 \node[tgTealVertex] (b) at (-.10,0) {$b$};
 \draw[tgEdge] (a)--(x)--(b);
 \draw[tgGold,-{Stealth[length=3mm]},line width=1.6pt] (.6,0)--(2.0,0);
 \node[tgBlueVertex] (c) at (2.8,0) {$a+1$};
 \node[tgTealVertex] (d) at (4.45,0) {$b+1$};
 \draw[tgAccent] (c)--(d);
 \node[font=\footnotesize,text=tgNavy] at (.55,-.97)
 {Usunięcie wewnętrznego wierzchołka $-1$ tworzy nową krawędź.};
\end{tikzpicture}
\end{center}

\begin{uwaga}[Gdy wynik nie jest zwykłym grafem]
Jeżeli usuwany wierzchołek o wadze $\pm1$ ma
kilku sąsiadów, eliminacja zmienia nie tylko ich
wagi, ale i wszystkie wzajemne elementy pozadiagonalne.
Może to prowadzić do krawędzi wielokrotnych,
ujemnych liczb splecenia lub niezerowych elementów
między wierzchołkami już sąsiadującymi. Dlatego
ogólny rachunek Kirby'ego wygodniej prowadzić
na \emph{ramowanych splotach} lub na macierzach
ze znakowanymi krawędziami, a nie na samych
nieskierowanych drzewach prostych.
\end{uwaga}

\section{Kontrola brzegu: wyznacznik i forma splecenia}
\label{sec:tg-kirby-linking}

\begin{stwierdzenie}[Postać formy splecenia po ruchu Kirby'ego]
Niech $Q$ będzie nieosobliwą macierzą splecenia
całkowicie ramowanego splotu i niech $Y$ będzie
brzegiem odpowiadającej mu czterorozmaitości uchwytowej.
Wtedy forma splecenia $\lambda_Y$ na
$H_1(Y;\mathbb Z)\cong\operatorname{coker}Q$ jest
opisywana macierzą $-Q^{-1}$ modulo $\mathbb Z$.
Przesunięcie uchwytu i kontrakcja składowej
$\varepsilon=\pm1$ zachowują klasę izometrii
tej formy, mimo że przy kontrakcji może zmienić
się ranga macierzy przecięcia.
\end{stwierdzenie}
\begin{proof}
Formułę na formę splecenia wyprowadziliśmy wcześniej
w sekcji \ref{sec:kt-linking-pairing}.
Nieosobliwość $Q$ zapewnia, że $Y$ jest wymierną
sferą homologiczną. Zmiana reprezentanta $x$
na $x+Qz$ zmienia $-x^TQ^{-1}y$
o liczbę całkowitą $-z^Ty$; analogicznie
dla $y$. Wzór jest więc dobrze określony. Przesunięcie
uchwytu daje $Q'=P^TQP$. Wówczas
${Q'}^{-1}=P^{-1}Q^{-1}P^{-T}$, a odwzorowanie
\[
 \Phi\colon \mathbb Z^r/Q\mathbb Z^r
 \longrightarrow \mathbb Z^r/Q'\mathbb Z^r,\qquad
 [x]\longmapsto[P^Tx]
\]
jest dobrze określone: $P^TQz=Q'P^{-1}z$,
gdzie $P^{-1}z\in\mathbb Z^r$.
Odwzorowanie odwrotne jest indukowane przez $P^{-T}$,
więc $\Phi$ jest izomorfizmem. Dla dowolnych
$x,y\in\mathbb Z^r$ mamy
\[
 -(P^Tx)^T{Q'}^{-1}(P^Ty)=-x^TQ^{-1}y.
\]
Po przejściu modulo $\mathbb Z$ otrzymujemy izometrię
form splecenia.

Przy kontrakcji otrzymujemy
$Q\cong[\varepsilon]\oplus B'$.
Z nieosobliwości $Q$ wynika nieosobliwość $B'$,
gdyż $\det Q=\varepsilon\det B'$.
Składnik $[\varepsilon]$ jest unimodularny i ma
trywialną grupę ilorazową. Odwrotność macierzy
blokowej to $[\varepsilon]\oplus(B')^{-1}$,
więc pierwszy składnik daje jedynie całkowite
wartości formy, znikające modulo $\mathbb Z$.
Stosując poprzedni rachunek izometrii do
całkowitej kongruencji, otrzymujemy dokładnie
formę odpowiadającą $B'$. Stąd w obu operacjach
forma brzegowa jest zachowana.
\end{proof}

\begin{przyklad}[Dwa równoważne diagramy i grupa $\mathbb Z/5$]
Niech $Q=\begin{psmallmatrix}-1&1\\1&-6\end{psmallmatrix}$.
Po usunięciu wierzchołka $-1$ dostajemy $B'=[-5]$.
Istotnie, $-6-(-1)^{-1}\cdot1^2=-5$.
Wyznaczniki to odpowiednio $5$ i $-5$,
więc ich moduły są jednakowe. Zakładamy tu standardowy diagram z węzłem
trywialnym o ramowaniu $-1$ oraz jego meridianem.
Obie chirurgie mają pierwszą grupę homologii
$\mathbb Z/5\mathbb Z$ (grupa ma rząd pierwszy),
a forma splecenia w każdym przypadku jest
przedstawiona, po odpowiedniej identyfikacji
generatora, przez $-[-5]^{-1}=1/5\pmod{\mathbb Z}$.
Ten przykład pokazuje, dlaczego wyznacznik
przy usuwaniu składnika $-1$ zmienia znak, podczas
gdy grupa brzegowa i jej forma pozostają niezmienione.
\end{przyklad}

\begin{przyklad}[Kontrakcja składowej $+1$ z dwoma sąsiadami]
\label{ex:tg-kirby-positive-blowdown}
Niech w ramowanym splocie występują trzy składowe:
węzeł trywialny $U$ o ramowaniu $+1$ i dwie
pozostałe składowe o ramowaniach
$-2$ i $-3$. Załóżmy, że
$\operatorname{lk}(U,K_1)=\operatorname{lk}(U,K_2)=1$,
a $\operatorname{lk}(K_1,K_2)=0$.
Macierz splecenia w kolejności $(U,K_1,K_2)$ wynosi
\[
Q=\begin{pmatrix}
1&1&1\\
1&-2&0\\
1&0&-3
\end{pmatrix}.
\]
Kontrakcja $U$ daje
\[
B'=
\begin{pmatrix}-2&0\\0&-3\end{pmatrix}
-\begin{pmatrix}1\\1\end{pmatrix}
\begin{pmatrix}1&1\end{pmatrix}
=\begin{pmatrix}-3&-1\\-1&-4\end{pmatrix},
\qquad \det B'=11.
\]
Zatem $\det Q=11$ i
$|H_1(\partial W_L;\mathbb Z)|=11$.
Ponieważ liczba $11$ jest pierwsza,
każda grupa abelowa tego rzędu jest cykliczna,
stąd $H_1(\partial W_L;\mathbb Z)\cong
\mathbb Z/11\mathbb Z$.
Sama obecność wpisu $1$ w macierzy
nie wystarczyłaby do wykazania cykliczności
grupy dla złożonego wyznacznika.
Na przykład macierz
$\operatorname{diag}(1,2,2)$ ma wpis $1$,
a jej kokernel jest izomorficzny z
$\mathbb Z/2\mathbb Z\oplus\mathbb Z/2\mathbb Z$.
Klasa pierwszego wektora bazowego $e_1$
w kokernelu $B'$ jest niezerowa:
$(B')^{-1}e_1=(-4,1)^T/11$ nie należy
do $\mathbb Z^2$. Ponieważ grupa ma rząd
pierwszy, klasa ta jest generatorem.
Na tym generatorze odpowiadającym $K_1$
forma splecenia ma wartość
\[
-\bigl((B')^{-1}\bigr)_{11}
=\frac4{11}\pmod{\mathbb Z}.
\]
Nowe składowe mają liczbę splecenia $-1$,
chociaż pierwotnie była ona zerowa.
Rachunek śledzi jednocześnie dodatnią
kontrakcję, ujemny znak nowej krawędzi,
grupę brzegową i postać splecenia.
\end{przyklad}

\section{Granice rachunku macierzowego i twierdzenie Kirby'ego}
\label{sec:tg-kirby-scope}

\begin{twierdzenie}[Kirby; postać klasyczna, bez dowodu]
Na mocy twierdzenia Lickorisha--Wallace'a każda zamknięta,
spójna, zorientowana trójrozmaitość powstaje przez
całkowitą chirurgię Dehna wzdłuż odpowiednio
ramowanego splotu w $S^3$.
Dwa takie przedstawienia dają dyfeomorficzne
trójrozmaitości z zachowaniem orientacji wtedy i tylko wtedy,
gdy można je połączyć skończoną liczbą izotopii
ramowanych splotów, przesunięć uchwytów oraz
dodawania lub usuwania składowej trywialnej
z ramowaniem $+1$ albo $-1$, \emph{odseparowanej}
od pozostałego splotu przez sferę $S^2$.
Sama rozłączność składowych, a nawet zerowe
algebraiczne liczby splecenia, nie wystarczają
do zastosowania tego ruchu.
\end{twierdzenie}

Twierdzenie o istnieniu przedstawienia chirurgicznego
przyjmujemy za D.~Rolfsenem, \emph{Knots and Links},
część 9I (por. wcześniejsze omówienie w rozdziale
o chirurgii Dehna). Zupełność ruchów Kirby'ego
jest zewnętrznym twierdzeniem:
R.~C. Kirby, \emph{A calculus for framed links
in $S^3$}, Invent. Math. 45 (1978), 35--56.
Przyjmujemy ją bez dowodu.
W kierunku zachowania brzegu izotopia przedłuża
się do dyfeomorfizmu otoczeń tubularnych,
przesunięcie uchwytu zachowuje brzeg na mocy
stwierdzenia \ref{prop:tg-kirby-congruence},
a dodanie odseparowanej składowej daje sumę spójną
brzegu z wynikiem chirurgii $\pm1$ na węźle
trywialnym. Ten ostatni wynik to $S^3$, jak
wykazano w dowodzie twierdzenia
\ref{thm:tg-kirby-blowdown}.
W kierunku przeciwnym, czyli od dyfeomorfizmu
brzegów do istnienia ciągu ruchów, stosujemy
właściwe twierdzenie Kirby'ego.

Ważne jest, że \emph{kongruencja macierzy}
nie jest kryterium równoważności chirurgii:
różne ramowane sploty mogą mieć tę samą macierz
splecenia, lecz dawać niehomeomorficzne brzegi. Na przykład
węzeł trywialny i lewoskrętny trójlistnik, oba
z ramowaniem $-1$, mają tę samą macierz $[-1]$.
Pierwsza chirurgia daje $S^3$, natomiast druga
daje sferę homologiczną Poincarégo
(por. R.~C. Kirby, M.~G. Scharlemann,
\emph{Eight faces of the Poincaré homology 3-sphere},
w: Geometric Topology, Academic Press, 1979,
opis 7, s.~121). Jej grupa fundamentalna
jest nietrywialna na mocy stwierdzenia
\ref{st:tg-E8-poincare}, więc brzegi nie są
nawet homeomorficzne bez uwzględniania orientacji.
Macierz rejestruje homologię i, gdy jest
nieosobliwa, formę splecenia na całej pierwszej
grupie homologii, a nie cały typ topologiczny
trójrozmaitości.

\begin{uwaga}[Granica między 3D i 4D]
Dodanie odseparowanego węzła trywialnego o ramowaniu
$+1$ lub $-1$ nie zmienia brzegu chirurgii,
ale na poziomie czterowymiarowym odpowiada
sumie spójnej we wnętrzu odpowiednio z
$\mathbb{CP}^2$ lub $\overline{\mathbb{CP}}^{\,2}$.
Istotnie, sfera separująca diagramy pozwala
rozpoznać ich czterowymiarowe wypełnienia jako
sumę brzegową wypełnień obu diagramów.
Wypełnienie pojedynczego węzła trywialnego
o ramowaniu $\pm1$ to odpowiednia zorientowana
wiązka dyskowa nad $S^2$, równoważnie
$\mathbb{CP}^2$ lub $\overline{\mathbb{CP}}^{\,2}$
z usuniętym wnętrzem kuli. Suma brzegowa
z takim składnikiem jest dyfeomorficzna
z opisaną sumą spójną we wnętrzu.
Przesunięcia uchwytów nie zmieniają samej
czterorozmaitości. Przy zastosowaniach do twierdzeń
Freedmana i Donaldsona trzeba więc zawsze pamiętać,
który rodzaj równoważności jest aktualnie rozważany.
\end{uwaga}

\begin{uwaga}[Połączenie z rozdziałem o $E_8$]
Grafy $A_n$, $D_n$ i $E_8$ są przykładami, w których
konstrukcje brzegu, jego homologii i formy przecięcia
można obliczyć bezpośrednio z macierzy grafu. Natomiast
identyfikacja konkretnej trójrozmaitości brzegowej
wymaga informacji geometrycznej: rozwarstwienia
Seiferta, grupy fundamentalnej, diagramu Kirby'ego
lub innych niezmienników. Dlatego oba nowe rozdziały
uzupełniają, ale nie zastępują, wcześniejszą teorię
węzłów i rozmaitości czterowymiarowych.
\end{uwaga}


% Iteracja 9. Nowy rozdział po rachunku Kirby'ego.
\begin{lemat}[Algebraiczna postać kontrakcji składowej o ramowaniu $\pm1$]
\label{lem:tg-audit-schur}
Niech $r\ge0$, $\varepsilon\in\{-1,1\}$,
$B\in M_r(\mathbb Z)$ będzie symetryczna,
$\ell\in\mathbb Z^r$ będzie kolumną i niech
$Q=\begin{psmallmatrix}B&\ell\\\ell^T&\varepsilon\end{psmallmatrix}$. Wówczas jest ona kongruentna nad
$\mathbb Z$ do
$\operatorname{diag}(B-\varepsilon\ell\ell^T,\varepsilon)$.
W szczególności
$\det Q=\varepsilon\det(B-\varepsilon\ell\ell^T)$.
\end{lemat}

\begin{proof}
Wystarczy wykonać całkowitą zmianę bazy, która zastępuje
$e_i$ przez $e_i-\varepsilon\ell_i e_{r+1}$ i zachowuje
$e_{r+1}$. W nowej bazie iloczyn
$(e_i-\varepsilon\ell_i e_{r+1})\cdot e_{r+1}$
wynosi $\ell_i-\varepsilon\ell_i\varepsilon=0$,
a iloczyn dwóch nowych wektorów jest równy
$B_{ij}-\varepsilon\ell_i\ell_j$.
Macierz przejścia jest trójkątna z wyznacznikiem jeden,
więc jest unimodularna. Daje to żądaną kongruencję
i, po wzięciu wyznaczników, wzór determinantowy.
Nie użyliśmy założenia nieosobliwości $B$ ani $Q$.
W szczególności dla $B'=B-\varepsilon\ell\ell^T$
kokernel $Q$ jest izomorficzny z kokernelem $B'$:
kongruencja całkowita daje izomorfizm grup
ilorazowych, a ostatni blok $[\varepsilon]$
ma trywialny kokernel. Ponadto
$\operatorname{rank}_{\mathbb Q}Q=
1+\operatorname{rank}_{\mathbb Q}B'$.
Przestrzenie jąder mają tę samą wymiarowość,
a sygnatury spełniają
$\operatorname{sign}Q=\operatorname{sign}B'+\varepsilon$.
Wszystkie te wnioski pozostają prawdziwe
dla $r=0$, przy konwencji wyznacznika pustej
macierzy równego jeden.
Dowodzi to wyłącznie części macierzowej operacji;
geometryczna realizacja kontrakcji jest osobnym krokiem
rachunku Kirby'ego.
\end{proof}

